\chapter{Il ML nei curricula della magistrale}\label{ch-02}

Questo pacco di slide è un ``intermezzo'' della prima lezione, che il
professore chiede di leggere autonomamente. Descrive come il corso di
Machine Learning si inserisce nella \textbf{Laurea Magistrale in
Informatica} dell'Università di Pisa e nei suoi curricula. Non contiene
materiale d'esame, ma aiuta a capire il ruolo del corso nel percorso di
studi.

\section{I curricula della
magistrale}\label{i-curricula-della-magistrale}

Dal 2017 la magistrale in Informatica è organizzata in quattro
curricula:

\begin{itemize}
\tightlist
\item
  \textbf{Artificial Intelligence} (AI);
\item
  \textbf{Data and Knowledge}, dal 2021 \textbf{Big Data Technologies}
  (BD);
\item
  \textbf{ICT Solutions Architect} (ICT);
\item
  \textbf{Software: Programming, Principles, and Technologies} (SW).
\end{itemize}

L'organizzazione per curricula permette di specializzarsi in un'area
(con un diploma supplementare che certifica il curriculum), offre corsi
metodologici di base all'inizio del percorso e mira a mostrare come due
anni di studio in più siano utili per il futuro professionale.
Regolamenti, istruzioni per gli esami e modalità di cambio di curriculum
o di piano di studi sono sul sito del dipartimento.

\section{Il ruolo del Machine
Learning}\label{il-ruolo-del-machine-learning}

Il ML è collegato a \textbf{tutti e quattro} i curricula, ma è un corso
\textbf{caratterizzante per il curriculum AI}, dove fa parte della base
metodologica per costruire sistemi adattivi e intelligenti insieme a:

\begin{itemize}
\tightlist
\item
  \emph{Artificial Intelligence Fundamentals} (6 CFU);
\item
  \emph{Computational Mathematics for Learning and Data Analysis} (9
  CFU);
\item
  \emph{Machine Learning} (9 CFU);
\item
  \emph{Parallel and Distributed Systems: Paradigms and Models} (9 CFU).
\end{itemize}

Gli altri caratterizzanti del curriculum AI sono \emph{Human Language
Technologies}, \emph{Intelligent Systems for Pattern Recognition} e
\emph{Smart Applications}, affiancati da gruppi di esami a scelta
(Algorithm Engineering, Data Mining, Mobile and Cyber-Physical Systems,
Information Retrieval, Computational Neuroscience, Robotics, Semantic
Web, ecc.).

Il ML è il perno dell'\textbf{area dei sistemi intelligenti}: molti
corsi della magistrale (Human Language Technologies, Intelligent Systems
for Pattern Recognition, Smart Applications, Computational Neuroscience,
Robotics, Data Mining, Information Retrieval, Bioinformatics, Big Data
Analytics\ldots) ne usano direttamente i metodi.

\begin{boxnote}{Piano di studi del curriculum AI (dal 2021)}

60 CFU di corsi caratterizzanti, 27 CFU da gruppi di esami a scelta (un
esame da 9 CFU e tre da 6 CFU) e almeno 9 CFU a libera scelta, che si
possono coprire anche con due esami da 6 CFU.

\end{boxnote}

\section{Sinergia con Computational
Mathematics}\label{sinergia-con-computational-mathematics}

Il professore consiglia di seguire in parallelo \textbf{ML} e
\textbf{CM} (\emph{Computational Mathematics for Learning and Data
Analysis}): il primo fornisce i metodi, il secondo le basi matematiche,
e i due corsi si stimolano a vicenda. CM non è comunque un prerequisito:
molti studenti di altri corsi di laurea o curricula hanno seguito ML
senza CM.

\section{Note per categorie particolari di
studenti}\label{note-per-categorie-particolari-di-studenti}

\begin{boxwarning}{Per chi viene dalla triennale di Pisa (corso IIA)}

Le prime sei lezioni sembreranno ``facili'', ma \textbf{non coincidono}
con i contenuti di IIA: contengono parti nuove e concetti matematici
diversi. Il cambio di prospettiva è netto: in IIA i modelli semplici
erano lo \emph{strumento} per introdurre i concetti di ML, qui i modelli
e i principi allo stato dell'arte sono \emph{l'oggetto} di studio. Dalla
lezione 5--6 in poi tutto è nuovo.

\end{boxwarning}

Per gli studenti del vecchio esame \textbf{AA1} (320AA, 6 CFU), che non
esiste più: chi deve passare da AA1 a ML (9 CFU) sostiene
un'integrazione su una lista di argomenti indicata dal professore; chi
deve fare il percorso inverso deve contattare il docente.

Per qualsiasi informazione il riferimento è il Prof.~Alessio Micheli
(micheli@di.unipi.it), del gruppo \emph{Computational Intelligence \&
Machine Learning} del Dipartimento di Informatica.
